% Generated by build_report.py. Edit the shared content there, then regenerate.
\documentclass[10pt,a4paper]{article}
\usepackage{fontspec}
\setmainfont{texgyretermes}[Extension=.otf,UprightFont=*-regular,BoldFont=*-bold,ItalicFont=*-italic,BoldItalicFont=*-bolditalic]
\usepackage[margin=1.9cm,headheight=14pt]{geometry}
\usepackage{amsmath,amssymb,booktabs,tabularx,longtable,array,ragged2e,enumitem,microtype,xcolor,tikz,fancyhdr,graphicx}
\usetikzlibrary{arrows.meta,positioning}
\usepackage[hidelinks,unicode]{hyperref}
\definecolor{ink}{HTML}{183A45}\definecolor{teal}{HTML}{087F7A}\definecolor{pale}{HTML}{EAF4F2}\definecolor{gray}{HTML}{56656B}
\newcolumntype{P}[1]{>{\RaggedRight\arraybackslash}p{#1}}
\setlength{\parindent}{0pt}\setlength{\parskip}{5pt}\setlength{\emergencystretch}{2em}
\setlength{\tabcolsep}{4pt}\renewcommand{\arraystretch}{1.16}
\setlist{nosep,leftmargin=1.4em,topsep=4pt}
\pagestyle{fancy}\fancyhf{}\fancyhead[L]{\small\color{gray}BẢO VỆ RIÊNG TƯ QUỸ ĐẠO}\fancyhead[R]{\small\color{gray}03/10/2026}\fancyfoot[C]{\small\thepage}\renewcommand{\headrulewidth}{0.3pt}
\newcommand{\key}[1]{\par\smallskip\noindent\colorbox{pale}{\parbox{\dimexpr\linewidth-2\fboxsep}{\textbf{#1}}}\par\smallskip}
\hypersetup{pdftitle={Bảo vệ riêng tư quỹ đạo: dữ liệu, phương pháp và kết quả},pdfauthor={}}
\begin{document}
\begin{center}{\LARGE\bfseries\color{ink}Bảo vệ riêng tư quỹ đạo}\\[5pt]{\large Related works, kiến trúc và kết quả đối chứng}\\[4pt]{\small Bản trình bày rút gọn · 03/10/2026 · Số liệu giữ từ bản 30/09}\end{center}
\section{Related works: mức bao phủ S1–S10}\label{sec:1}
Khảo sát giữ cửa sổ \textbf{21/09/2023–21/09/2026}. Bảng đối chiếu từng paper với mười nhiệm vụ của ta; mức bao phủ phải đọc cùng giả định và giới hạn ở cột cuối.

S1: vị trí hiện tại; S2: nơi dừng; S3: đoạn đường đã đi; S4: liên kết danh tính giữa phiên; S5: cạnh đường kế tiếp; S6: đích chưa tới; S7: nội dung/ý định truy vấn; S8: định vị nhờ dữ liệu người đi cùng; S9: điểm đầu bị che; S10: điểm cuối bị che của chuyến đã hoàn tất.

\textbf{Fully (F):} paper trực tiếp xử lý mục tiêu suy luận của scenario, trong giả định của paper. \textbf{Partially (P):} chỉ xử lý một phần hoặc bối cảnh hẹp hơn. \textbf{CX:} chưa đủ nguồn để phân loại. \textbf{—:} chưa tìm thấy bằng chứng xử lý trực tiếp; không kết luận cơ chế chắc chắn thất bại. F/P là đối chiếu của ta, không phải nhãn do tác giả công bố, không có nghĩa bảo vệ 100\% hoặc đã vượt các ca A/B/C của ta.

\begingroup\small
\begin{longtable}{@{}P{3.387cm}P{1.742cm}P{2.613cm}P{8.614cm}@{}}
\toprule
\textbf{Paper} & \textbf{Fully} & \textbf{Partially / CX} & \textbf{Căn cứ và giới hạn}\\\midrule
\endfirsthead
\toprule
\textbf{Paper} & \textbf{Fully} & \textbf{Partially / CX} & \textbf{Căn cứ và giới hạn}\\\midrule
\endhead
\bottomrule\endfoot
TransProtect 2024 [\hyperlink{ref-R1}{R1}] & S1, S3 & S2 & §3–5: suy vị trí từ chuỗi đã làm nhiễu, biết đường/traffic; chưa đánh giá riêng nơi dừng.\\
\addlinespace[3pt]
Semantic correlation 2026 [\hyperlink{ref-R2}{R2}] & S1, S3 & S2 & §3–6: giữ dummy hợp lý theo thời gian/ngữ nghĩa; chưa chấm riêng nơi dừng hay ý định S7.\\
\addlinespace[3pt]
Fake queries 2026 [\hyperlink{ref-R3}{R3}] & S1, S3 & S2 & §3–6: chèn truy vấn giả để khó nối tuyến; chưa đánh giá riêng tọa độ nơi dừng.\\
\addlinespace[3pt]
Road-aware PTPPM 2025 [\hyperlink{ref-R4}{R4}] & S1, S3 & S2 & §III, VI: chống suy vị trí/quỹ đạo có tương quan; bước dự báo chưa là phép thử S5/S6.\\
\addlinespace[3pt]
CPCROK 2025 [\hyperlink{ref-R5}{R5}] & — & S4 & §3–5: chống nối bí danh xe trong VANET; chưa tách danh tính người/thiết bị như S4.A/B/C.\\
\addlinespace[3pt]
DP-FETC 2025 [\hyperlink{ref-R6}{R6}] & — & S3, S8, S9, S10 & SynFE/PreOD: sinh dữ liệu offline, đổi đầu/cuối của tuyến tương quan; không chấm định vị từ người đi cùng.\\
\addlinespace[3pt]
Improved PIR 2025 [\hyperlink{ref-R7}{R7}] & CX & CX: S1, S7 & Tóm tắt nêu bảo vệ vị trí/nội dung bằng mã hóa; thiếu toàn văn để đối chiếu quyền quan sát và suy ý định.\\
\addlinespace[3pt]
SoK 2024 [\hyperlink{ref-R8}{R8}] & Không áp dụng & Không áp dụng & Khảo sát cách đánh giá/tấn công; không phải cơ chế bảo vệ để gán F/P.\\
\addlinespace[3pt]
PRISM 2025 [\hyperlink{ref-R9}{R9}] & CX & CX: S1, S3 & Tóm tắt: bảo vệ chia sẻ quỹ đạo; chưa đủ mô hình attacker. Dùng LSTM không xác nhận xử lý S5/S6.\\
\addlinespace[3pt]
Options to Action 2026 [\hyperlink{ref-R10}{R10}] & Không áp dụng & Không áp dụng & Khảo sát sử dụng tính năng, có che đầu/cuối (S9/S10); không đo khả năng chống suy luận.\\
\addlinespace[3pt]
Forsch et al. 2023 [\hyperlink{ref-R11}{R11}] & S9, S10 & — & §3.2: S-TT cắt cả hai đầu, xét nhiều tuyến cùng đích; bảo đảm theo tập tấn công/địa điểm cho trước, offline.\\
\addlinespace[3pt]
EV querying 2024 [\hyperlink{ref-R12}{R12}] & S1, S3 & S2 & Proposed Query Model: AGeoI + dummy chống định vị/nối tuyến online; chưa chấm riêng nơi dừng hoặc S4–S10.\\
\end{longtable}
\endgroup
Scenario không được liệt kê ở một hàng được hiểu là —; với R7/R9 là chưa xác minh, với R8/R10 là không áp dụng. Bảng cho thấy các cơ chế chuỗi chủ yếu xét S1/S3; S2 cần phép thử nơi dừng riêng, còn S9/S10 cần kiểm tra suy ngược từ phần tuyến vẫn công bố.


\clearpage
\subsection{Metrics gốc và điều rút ra từ related works}\label{sub:1.1}
Bảng dưới ghi các chỉ số đã xác minh trong nguồn. $\uparrow$/$\downarrow$ là hướng tốt hơn theo mục tiêu của paper; CX nghĩa là chưa xác minh đủ định nghĩa. Các paper dùng dataset và nhiệm vụ khác nhau, nên không so trực tiếp điểm số gốc của chúng.

\begingroup\small
\begin{longtable}{@{}P{3.871cm}P{4.355cm}P{3.871cm}P{4.258cm}@{}}
\toprule
\textbf{Paper / hướng nghiên cứu} & \textbf{Privacy / chẩn đoán} & \textbf{Utility} & \textbf{Chi phí / phạm vi}\\\midrule
\endfirsthead
\toprule
\textbf{Paper / hướng nghiên cứu} & \textbf{Privacy / chẩn đoán} & \textbf{Utility} & \textbf{Chi phí / phạm vi}\\\midrule
\endhead
\bottomrule\endfoot
TransProtect 2024 [\hyperlink{ref-R1}{R1}] & EIE $\uparrow$: sai số suy luận & $\Delta$c $\downarrow$: chênh chi phí hành trình & Byte + độ trễ toàn dịch vụ: CX\\
\addlinespace[3pt]
Semantic correlation 2026 [\hyperlink{ref-R2}{R2}] & ASR $\uparrow$: đạt ẩn danh; DER $\uparrow$: dummy/ngữ nghĩa & Chưa đo top-5 POI bằng DER & Thời gian sinh $\downarrow$\\
\addlinespace[3pt]
Fake queries 2026 [\hyperlink{ref-R3}{R3}] & Số đường khó phân biệt $\uparrow$; ASR $\uparrow$ & Chất lượng POI: CX & Thời gian sinh, RSS $\downarrow$\\
\addlinespace[3pt]
Road-aware PTPPM 2025 [\hyperlink{ref-R4}{R4}] & Sai số $\uparrow$; tấn công đúng $\downarrow$ & Dịch chuyển đầu ra $\downarrow$ & Chi phí LBS đầu-cuối: CX; preprint\\
\addlinespace[3pt]
CPCROK 2025 [\hyperlink{ref-R5}{R5}] & Tỷ lệ khớp quỹ đạo $\rho$ $\downarrow$ & Bối cảnh VANET & Số quỹ đạo giả $\psi$ $\downarrow$; ước lượng gián tiếp chi phí\\
\addlinespace[3pt]
DP-FETC 2025 [\hyperlink{ref-R6}{R6}] & MI $\downarrow$; bảo đảm $\varepsilon$-DP & JSD phân bố $\downarrow$ & Độ phức tạp; xuất bản offline\\
\addlinespace[3pt]
Improved PIR 2025 [\hyperlink{ref-R7}{R7}] & Phân tích an toàn; metric tấn công: CX & Truy hồi cụ thể: CX & Chi phí tính toán; mới đủ tóm tắt\\
\addlinespace[3pt]
SoK 2024 [\hyperlink{ref-R8}{R8}] & Tấn công tái dựng/liên kết & Phân bố, hình học, tác vụ & Khảo sát, không phải model\\
\addlinespace[3pt]
PRISM 2025 [\hyperlink{ref-R9}{R9}] & LPI; công thức CX & TDD, DC; định nghĩa CX & Thời gian xử lý/phản hồi\\
\addlinespace[3pt]
Options to Action 2026 [\hyperlink{ref-R10}{R10}] & Tỷ lệ dùng tính năng & Nhận thức và sử dụng & Nghiên cứu hành vi\\
\addlinespace[3pt]
Forsch et al. 2023 [\hyperlink{ref-R11}{R11}] & Trình bày S-TT / k-site-unidentifiability & Che điểm nhạy cảm & Chương tổng quan; kế thừa S-TT 2022\\
\addlinespace[3pt]
EV querying 2024 [\hyperlink{ref-R12}{R12}] & AGeoI + dummy & Dịch vụ truy vấn trạm sạc & Chưa xác minh phép thử điểm đầu/cuối\\
\end{longtable}
\endgroup
\textbf{Mốc nền giữ riêng:} DLS 2014 [\hyperlink{ref-B1}{B1}], RDG 2021 [\hyperlink{ref-B2}{B2}]; AnotherMe [\hyperlink{ref-B3}{B3}] có số tạp chí 2024 nhưng online 11/09/2023, ngoài cửa sổ. S-TT 2022 [\hyperlink{ref-B4}{B4}], tấn công EPZ 2022 [\hyperlink{ref-B5}{B5}] và Data Stream Obfuscator 07/2023 [\hyperlink{ref-B6}{B6}] bổ sung góc nhìn điểm đầu/cuối. Không gọi các nguồn này là paper mới trong ba năm.

\key{Các paper đo những khía cạnh khác nhau. Để so trên bài toán của ta, cần cùng đo khả năng suy luận của đối thủ, chất lượng dịch vụ và chi phí.}


\clearpage
\subsection{Vì sao chọn bộ metrics đề xuất?}\label{sub:1.2}
\begingroup\small
\begin{longtable}{@{}P{5.513cm}P{4.135cm}P{6.990cm}@{}}
\toprule
\textbf{Vấn đề từ related works} & \textbf{Metric đề xuất} & \textbf{Ý nghĩa trong bài toán của ta}\\\midrule
\endfirsthead
\toprule
\textbf{Vấn đề từ related works} & \textbf{Metric đề xuất} & \textbf{Ý nghĩa trong bài toán của ta}\\\midrule
\endhead
\bottomrule\endfoot
Độ bất định, tương đồng dummy và lượng thông tin rò rỉ đo các khía cạnh khác nhau. ASR cũng không có cùng định nghĩa giữa các paper. & Hit50/100/200 $\downarrow$; MAE, median, p90 $\uparrow$ & Đối thủ (attacker) phải đoán tọa độ thật từ dữ liệu công bố. Hit100 là tỷ lệ đoán cách đáp án $\le$100 m, dùng được cho cả tập dummy và vị trí thay thế.\\
\addlinespace[3pt]
MAE lớn vẫn có thể che một số lần suy đúng rất gần. & Đọc Hit cùng phân bố sai số & MAE là sai số trung bình; median là trung vị; p90 là ngưỡng chứa 90\% sai số mẫu. S1 chấm một điểm, S2 nơi dừng, S3 chuỗi điểm, S9/S10 điểm đầu/cuối.\\
\addlinespace[3pt]
Bảo toàn phân bố hoặc giảm độ lệch chưa cho biết người dùng có nhận đúng địa điểm cần tìm. & Recall@5 $\uparrow$; trả đủ $\uparrow$; $\Delta$D đường $\downarrow$ & So 5 địa điểm được chọn tại thiết bị với 5 địa điểm chuẩn theo GPS thật, cùng trạng thái máy chủ. Mỗi ca cần Recall $\ge$90\%.\\
\addlinespace[3pt]
Chỉ đếm dummy hoặc thời gian sinh sẽ bỏ sót phản hồi, truy vấn chèn và tải ngoài chuyến. & Byte yêu cầu + phản hồi $\downarrow$; latency p50/p95 $\downarrow$ & Đếm toàn bộ dữ liệu gửi/nhận cho mỗi lần dùng dịch vụ thật. Độ trễ toàn dịch vụ chưa đo; chưa có điểm tổng Q thực nghiệm.\\
\end{longtable}
\endgroup
POI là địa điểm như nhà thuốc hoặc trạm sạc. Recall@5 = số POI chuẩn tìm được / số POI chuẩn (tối đa 5). Ví dụ, tìm đúng 4/5 địa điểm được 80\%. Có đáp án mà không phục vụ thì tính 0; không có đáp án chuẩn thì để thiếu. “Trả đủ” đo số lượng; $\Delta$D đo khoảng cách đường tăng thêm.

Để so công bằng, giữ cùng mục tiêu, quyền quan sát và dịch vụ: máy chủ trả top-10, thiết bị chọn top-5. Với mỗi loại dữ liệu công bố (transcript), chọn bộ suy luận phù hợp trên nhóm phát triển rồi giữ cố định khi kiểm tra. Gộp các lần chạy trong bản ghi và nhóm tuyến; cho các ca trọng số bằng nhau trong scenario, rồi cho năm scenario trọng số bằng nhau. S10 chỉ có A/B. Chi phí truyền vẫn phải báo riêng.

\key{Bộ đo trả lời ba câu hỏi: đối thủ đoán đúng tới đâu, người dùng nhận đúng kết quả đến đâu và cần bao nhiêu chi phí. Đóng góp là cách chọn và áp dụng chung các chỉ số đã có cho nhiệm vụ này.}


\clearpage
\section{Kiến trúc mô hình BR-Dummy: Geo-I, mạng đường và các cơ chế}\label{sec:2}
\textbf{Nền tảng là Geo-I:} từ GPS thật, mô hình lấy ngẫu nhiên một vị trí đại diện đã được bảo vệ, gọi là \textbf{neo}. Lịch sử neo, mạng đường và POI công khai được dùng để sinh K vị trí/quỹ đạo giả. Lớp xử lý biên (boundary) bổ sung bảo vệ điểm đầu/cuối.

\begin{center}
\includegraphics[width=\linewidth]{figures/architecture_model.pdf}
\end{center}
Đọc từ trái sang phải: input $\to$ layer bảo vệ Geo-I $\to$ layer ngữ cảnh $\to$ layer sinh dummy $\to$ output. Khung xanh bao toàn bộ mô hình của ta, gồm ba layer lõi và layer 4 boundary S9/S10 ở hàng dưới; input và bên nhận ở ngoài khung. Switching, pacing và slack tùy biến thể.

\begingroup\small
\begin{longtable}{@{}P{4.105cm}P{12.814cm}@{}}
\toprule
\textbf{Component} & \textbf{Cơ chế và vai trò đối với scenario}\\\midrule
\endfirsthead
\toprule
\textbf{Component} & \textbf{Cơ chế và vai trò đối với scenario}\\\midrule
\endhead
\bottomrule\endfoot
Geo-I / neo REM & Lấy neo bằng cơ chế mũ REM: vị trí càng xa GPS theo khoảng cách Euclid, xác suất được chọn càng thấp. Tập ứng viên đường được xác định công khai; không cắt ứng viên theo bán kính riêng quanh GPS. Đây là nền bảo vệ S1.\\
\addlinespace[3pt]
Tái dùng neo + ngân sách & Kiểm tra khoảng cách có nhiễu để giữ neo cũ hoặc tạo neo mới; cả hai bước đều tiêu ngân sách riêng tư. Tái dùng neo giúp giảm số mẫu độc lập khi dừng (S2). Ngân sách giới hạn rò rỉ tích lũy qua chuỗi (S3); pacing giãn các lần đọc GPS.\\
\addlinespace[3pt]
Belief từ lịch sử neo & Belief là phân bố ước lượng xe có thể đang ở đâu, được cập nhật từ neo đã bảo vệ. Biến thể switching phân biệt trạng thái dừng/đi. Khối này hỗ trợ chọn dummy cho S2/S3, không dùng tốc độ thật hay đích tương lai.\\
\addlinespace[3pt]
Road-aware + POI-aware & Road-aware giới hạn dummy ở nơi có thể tới theo hướng làn và luật rẽ. POI-aware chọn K điểm để phủ nhiều POI bổ sung nhau. Exchange đổi ứng viên; slack cho phép độ lệch khi chọn, tùy biến thể. Hỗ trợ chuỗi hợp lý (S3) và dịch vụ; sức chống suy luận vẫn cần đo.\\
\addlinespace[3pt]
S9: bỏ đầu trước lõi & Không đưa GPS của h giây đầu vào lõi. Vì vậy, những điểm đầu bị che không đi vào trạng thái neo rồi ảnh hưởng các bản tin sau. Đối thủ vẫn có thể suy ngược từ tuyến còn thấy.\\
\addlinespace[3pt]
S10: buffer sau lõi & Giữ mỗi bản tin đã bảo vệ trong bộ đệm $\Delta$ giây rồi mới phát. Khi chuyến kết thúc, hủy bản tin chưa phát. Cách này che phần cuối mà không cần biết trước đích, nhưng làm tăng độ trễ và chưa loại hết suy luận từ phần còn thấy.\\
\end{longtable}
\endgroup
Belief, mạng đường và bộ chọn POI chỉ xử lý neo đã bảo vệ cùng dữ liệu công khai. Theo tính chất hậu xử lý, chúng giữ bảo đảm của neo khi tính đúng ngân sách qua các lần truy cập (composition). K dummy không tự tạo k-anonymity. GPS thật chỉ thêm ở bước chọn top-5 tại thiết bị sau phản hồi, không đưa vào bộ chọn dummy. Pacing chưa che giờ chuyến. Boundary đã kiểm tra mã/tích hợp; chưa có benchmark toàn mô hình với dịch vụ.


\clearpage
\section{Benchmark: phân biệt lõi Geo-I và nhánh truy vấn công khai}\label{sec:3}
\key{Phạm vi bảng dưới: Đề xuất 30/67 là nhánh kế hoạch truy vấn công khai, không dùng neo Geo-I. Đây là kết quả của nhánh mở rộng dịch vụ; không gán Hit100=0 cho kiến trúc Geo-I/BR-Dummy hoặc BR-Boundary ở mục 2.}

Các đối chứng chạy cùng dịch vụ và được chấm trên cùng mục tiêu. \textbf{Năm bản triển khai điều chỉnh từ paper (adapter), chạy trực tuyến:} DLS, RDG, TransProtect*, Semantic*, Fake-query*. * TransProtect dùng Markov thay Transformer; Semantic dùng bộ dự báo thực nghiệm thay LSTM; Fake-query dùng giả định lịch chèn công khai. AnotherMe là tham chiếu offline, báo riêng phía dưới.

\subsection{Nhánh công khai: S1–S3 trên bốn nhóm tuyến}\label{sub:3.1}
54 bản ghi từ 30 chuyến trong bốn nhóm 901–904; đối thủ được huấn luyện/chọn trên nhóm khác. Bảng là \textbf{Hit100 $\downarrow$}, trung bình đều A/B/C; chất lượng dịch vụ tại p=0,8 (80\% POI khả dụng), trung bình đều năm scenario. Đây chưa phải kết quả của cơ chế gửi theo lịch trên bộ 32 nhóm.

\begingroup\small
\begin{longtable}{@{}P{3.425cm}P{1.998cm}P{1.998cm}P{1.998cm}P{2.854cm}P{3.520cm}@{}}
\toprule
\textbf{Phương pháp} & \textbf{S1} & \textbf{S2} & \textbf{S3} & \textbf{Recall@5 $\uparrow$} & \textbf{Byte / sự kiện $\downarrow$}\\\midrule
\endfirsthead
\toprule
\textbf{Phương pháp} & \textbf{S1} & \textbf{S2} & \textbf{S3} & \textbf{Recall@5 $\uparrow$} & \textbf{Byte / sự kiện $\downarrow$}\\\midrule
\endhead
\bottomrule\endfoot
Vị trí thật & 100,00\% & 100,00\% & 100,00\% & 100,00\% & 242,8\\
\addlinespace[3pt]
DLS & 50,00\% & 41,67\% & 95,76\% & 100,00\% & 962,7\\
\addlinespace[3pt]
RDG & 58,33\% & 62,50\% & 95,83\% & 100,00\% & 950,7\\
\addlinespace[3pt]
TransProtect* & 66,67\% & 79,17\% & 74,59\% & 99,69\% & 242,7\\
\addlinespace[3pt]
Semantic* & 83,33\% & 100,00\% & 84,39\% & 100,00\% & 967,4\\
\addlinespace[3pt]
Fake-query* & 50,00\% & 33,33\% & 96,91\% & 100,00\% & 1.313,2\\
\addlinespace[3pt]
Đề xuất 30 & 0,00\% & 0,00\% & 0,00\% & 94,75\% & 2.389,7\\
\addlinespace[3pt]
Đề xuất 67 & 0,00\% & 0,00\% & 0,00\% & 100,00\% & 5.353,9\\
\end{longtable}
\endgroup
\textbf{Cách đọc:} vị trí thật và các adapter vẫn cho đối thủ suy đúng gần ở một phần ca; hai kế hoạch công khai đạt Hit100=0 trong tập đối thủ đã thử. Điều này phù hợp với việc tọa độ công bố không bám GPS, nơi dừng hay tuyến thật. Không diễn giải 0\% thành an toàn trước mọi đối thủ.

\textbf{Chất lượng dịch vụ và giá phải trả:} bản 30 giảm traffic so với 67 nhưng độ phủ POI thấp hơn. Bản 67 đạt Recall cao hơn nhờ thêm query; cần đọc riêng cả ba mặt: riêng tư, dịch vụ và chi phí. Byte là request + response JSON, giữ cả query phụ; chưa gồm HTTP/TLS. Bản 67 còn tám POI ngoài độ phủ tĩnh, nên Recall=100\% trên mẫu chưa là chứng nhận toàn bản đồ.


\clearpage
\subsection{Nhánh công khai: S9–S10 trên 32 nhóm tuyến, gửi theo lịch công khai}\label{sub:3.2}
Bộ 32 nhóm có 704 chuyến. Chất lượng dịch vụ được chấm trên 435 bản ghi từ 246 chuyến; riêng tư điểm đầu/cuối trên 150 bản ghi. Đối thủ được chọn cố định trước khi kiểm tra bộ này, có thể suy xuôi/ngược theo mạng đường và kết hợp chuyến lặp. S9 gộp A/B/C, S10 gộp A/B; mỗi ô privacy là \textbf{Hit100 $\downarrow$ / MAE (m) $\uparrow$}.

\begingroup\small
\begin{longtable}{@{}P{3.177cm}P{3.658cm}P{3.658cm}P{2.695cm}P{2.888cm}@{}}
\toprule
\textbf{Phương pháp} & \textbf{S9: Hit / MAE} & \textbf{S10: Hit / MAE} & \textbf{Recall@5 $\uparrow$} & \textbf{Byte / sự kiện $\downarrow$}\\\midrule
\endfirsthead
\toprule
\textbf{Phương pháp} & \textbf{S9: Hit / MAE} & \textbf{S10: Hit / MAE} & \textbf{Recall@5 $\uparrow$} & \textbf{Byte / sự kiện $\downarrow$}\\\midrule
\endhead
\bottomrule\endfoot
Vị trí thật & 42,48\% / 165 & 17,39\% / 212 & 100,00\% & 249,9\\
\addlinespace[3pt]
DLS & 34,33\% / 462 & 11,16\% / 222 & 100,00\% & 952,2\\
\addlinespace[3pt]
RDG & 19,76\% / 793 & 10,75\% / 224 & 100,00\% & 944,4\\
\addlinespace[3pt]
TransProtect* & 33,61\% / 243 & 13,31\% / 259 & 99,70\% & 249,9\\
\addlinespace[3pt]
Semantic* & 40,81\% / 184 & 11,82\% / 292 & 100,00\% & 975,8\\
\addlinespace[3pt]
Fake-query* & 31,36\% / 522 & 26,11\% / 213 & 100,00\% & 1.295,6\\
\addlinespace[3pt]
30 + lịch & 0,00\% / 2.232 & 0,00\% / 1.916 & 95,00\% & 3.071,0\\
\addlinespace[3pt]
67 + lịch & 0,00\% / 2.311 & 0,00\% / 1.916 & 100,00\% & 6.878,6\\
\end{longtable}
\endgroup
\textbf{Privacy:} lịch công khai đạt Hit100=0 so với S9 \textbf{19,76–40,81\%} và S10 \textbf{10,75–26,11\%} của năm adapter. Đối chứng dùng GPS thật vẫn cho phép suy điểm cuối ở cả S10.A (12,90\%) và S10.B (21,88\%), nên phép thử vẫn phát hiện được rò rỉ khi có. Khoảng tin cậy (CI) 95\% của chênh lệch Hit100 giữa đề xuất và từng adapter nằm dưới 0 ở S9/S10, hỗ trợ kết luận đề xuất có tỷ lệ đoán đúng thấp hơn. Tính bằng 3.000 lần lấy mẫu lại nhóm tuyến (bootstrap); chưa hiệu chỉnh nhiều so sánh.

\textbf{Chất lượng dịch vụ:} tại p=0,8, bản 30 đạt 95,00\%, 14/14 ca đạt $\ge$90\%; bản 67 đạt 100,00\%, 14/14 ca. Khi tăng tỷ lệ POI khả dụng lên p=0,95: bản 30 đạt 93,01\%, 13/14 ca; bản 67 đạt 100\%, 14/14 ca ở các mức đã thử.

\textbf{Chi phí:} khoảng 3.071 / 6.879 byte mỗi sự kiện thật. So với cùng kế hoạch chỉ gửi lại truy vấn khi có hoạt động, lịch công khai tốn 4,48 lần byte: đây là giá của việc che giờ hoạt động, không phải tăng độ phủ POI.

\textbf{AnotherMe:} paper mô tả hệ thống online, nhưng adapter VTGA hiện có đọc cả chuyến. Trên 32 nhóm, riêng tư chỉ chấm trên tập con sinh đầu ra thành công có Hit100=0, Recall toàn phạm vi chỉ 16,85\%. Ca lỗi/thiếu đầu ra vẫn tính chất lượng dịch vụ bằng 0, chỉ số riêng tư để thiếu; không so trực tiếp các số này như có cùng mẫu số với năm adapter online.

\key{Kết luận cho nhánh công khai: thiết kế phối hợp kế hoạch công khai, lịch tải và xếp hạng cục bộ có tỷ lệ suy đúng thấp hơn trong năm scenario đã thử, đổi lại traffic cao hơn. Chưa chứng minh vượt các paper nguyên bản hoặc tốt hơn ở cùng ngân sách; dữ liệu vẫn cùng thành phố/bộ sinh SUMO.}


\clearpage
\section{Phụ lục: đọc chi tiết 14 sample A/B/C}\label{sec:4}
A/B/C là các điều kiện cùng nhiệm vụ, không phải thứ tự độ khó. Mỗi ca chọn một bản ghi thật từ bộ minh họa 12 nhóm/264 chuyến; số bản ghi không phải số chuyến độc lập. \textbf{S10 chỉ có A/B} trong tài liệu; B ánh xạ mã nguồn cũ S10.C để giữ truy nguyên, không có ca C độc lập trong phạm vi hiện tại.

Chấm màu = mẫu trước bảo vệ; nét đứt = đường thật để chấm; sao đỏ = đáp án; trục thời gian tách các mẫu chồng nhau. Đối thủ chỉ nhận dữ liệu đã bảo vệ và thông tin phụ trợ được phép, không nhận các tọa độ/nhãn thật này. Bản đồ minh họa dùng nền SUMO với cùng mã kiểm tra tệp OSM; thiếu mạng gốc để xác minh trùng toàn bộ hình học. © OpenStreetMap contributors.

\subsection{S1 — suy ra vị trí tại một lần gửi}\label{sub:4.1}
Đối thủ cần suy ra tọa độ tại một lần gửi. A/B thay ràng buộc mạng đường; C thêm ngữ cảnh POI hiếm.

\noindent\begin{minipage}[t]{0.48\linewidth}\vspace{0pt}\includegraphics[width=\linewidth]{figures/case_S1_A.pdf}\end{minipage}\hfill\begin{minipage}[t]{0.49\linewidth}\vspace{0pt}\RaggedRight \textbf{S1.A / 12 bản ghi.} Một bản tin tại cạnh có $\ge$2 hướng đi tiếp hợp lệ: kiểm tra khi mạng đường còn nhiều khả năng.\par\smallskip \textbf{Bản ghi minh họa v2-r301-000.} u301\_\allowbreak{}00: 1 mẫu, chỉ số GPS FCD[333…333]\par\smallskip\textbf{Đọc sample.} Một lần lấy mẫu ở cạnh có 4 hướng đi tiếp. Sao đỏ trùng vị trí đầu vào; đối thủ phải chọn vị trí từ đầu ra đã bảo vệ, không được thấy sẵn sao này.\end{minipage}\par\vspace{3pt}
\noindent\begin{minipage}[t]{0.48\linewidth}\vspace{0pt}\includegraphics[width=\linewidth]{figures/case_S1_B.pdf}\end{minipage}\hfill\begin{minipage}[t]{0.49\linewidth}\vspace{0pt}\RaggedRight \textbf{S1.B / 12 bản ghi.} Một bản tin tại cạnh chỉ có 1 hướng đi tiếp: bản đồ tự loại bớt khả năng.\par\smallskip \textbf{Bản ghi minh họa v2-r301-001.} u301\_\allowbreak{}00: 1 mẫu, chỉ số GPS FCD[241…241]\par\smallskip\textbf{Đọc sample.} Vẫn chỉ một lần lấy mẫu, nhưng cạnh chỉ có một hướng đi tiếp. Bản đồ có thể loại bớt khả năng so với A; cần đo lợi thế này bằng đối thủ thực tế.\end{minipage}\par\vspace{3pt}
\noindent\begin{minipage}[t]{0.48\linewidth}\vspace{0pt}\includegraphics[width=\linewidth]{figures/case_S1_C.pdf}\end{minipage}\hfill\begin{minipage}[t]{0.49\linewidth}\vspace{0pt}\RaggedRight \textbf{S1.C / 12 bản ghi.} Cách POI $\le$150 m; loại POI chiếm $\le$5\% danh mục công khai. Hiếm theo loại địa điểm, không phải theo tần suất ghé thăm.\par\smallskip \textbf{Bản ghi minh họa v2-r301-029.} u301\_\allowbreak{}13: 1 mẫu, chỉ số GPS FCD[695…695]\par\smallskip\textbf{Đọc sample.} Điểm lấy mẫu cách POI pharmacy 12,03 m; loại này chiếm 3,11\% danh mục. Ngữ cảnh địa điểm hiếm là dấu hiệu bổ sung, không phải chuỗi di chuyển.\end{minipage}\par\vspace{3pt}

\clearpage
\subsection{S2 — suy ra nơi dừng qua nhiều lần gửi}\label{sub:4.2}
Đối thủ cần suy ra tọa độ nơi dừng. A/B thay thời lượng và nhịp quan sát; C kiểm tra rời đi rồi quay lại.

\noindent\begin{minipage}[t]{0.48\linewidth}\vspace{0pt}\includegraphics[width=\linewidth]{figures/case_S2_A.pdf}\end{minipage}\hfill\begin{minipage}[t]{0.49\linewidth}\vspace{0pt}\RaggedRight \textbf{S2.A / 12 bản ghi.} Dừng có chủ đích $\ge$20 giây; lấy truy vấn mỗi 5 giây.\par\smallskip \textbf{Bản ghi minh họa v2-r301-002.} u301\_\allowbreak{}05: 6 mẫu, chỉ số GPS FCD[54…79]\par\smallskip\textbf{Đọc sample.} 6 mẫu cùng một tọa độ, cách nhau 5 giây; đoạn dừng thật dài 29 giây. Sáu vạch thời gian bị chồng thành một chấm trên map.\end{minipage}\par\vspace{3pt}
\noindent\begin{minipage}[t]{0.48\linewidth}\vspace{0pt}\includegraphics[width=\linewidth]{figures/case_S2_B.pdf}\end{minipage}\hfill\begin{minipage}[t]{0.49\linewidth}\vspace{0pt}\RaggedRight \textbf{S2.B / 12 bản ghi.} Dừng $\ge$120 giây; lấy truy vấn mỗi 20 giây.\par\smallskip \textbf{Bản ghi minh họa v2-r301-003.} u301\_\allowbreak{}06: 9 mẫu, chỉ số GPS FCD[59…219]\par\smallskip\textbf{Đọc sample.} 9 mẫu cùng một tọa độ, cách nhau 20 giây; đoạn dừng dài 179 giây. Đối thủ có cửa sổ dài hơn để kết hợp bằng chứng, không chỉ một bản tin.\end{minipage}\par\vspace{3pt}
\noindent\begin{minipage}[t]{0.48\linewidth}\vspace{0pt}\includegraphics[width=\linewidth]{figures/case_S2_C.pdf}\end{minipage}\hfill\begin{minipage}[t]{0.49\linewidth}\vspace{0pt}\RaggedRight \textbf{S2.C / 12 bản ghi.} Hai lần dừng trên cùng làn, mỗi lần $\ge$20 giây; ở giữa thực sự di chuyển.\par\smallskip \textbf{Bản ghi minh họa v2-r301-004.} u301\_\allowbreak{}07: 18 mẫu, chỉ số GPS FCD[214…1193]\par\smallskip\textbf{Đọc sample.} 18 mẫu chia hai đợt, mỗi đợt 9 mẫu. Hai lần dừng dài 44 giây tại cùng tọa độ; khoảng trống giữa hai cụm thời gian là lúc xe rời đi rồi quay lại.\end{minipage}\par\vspace{3pt}

\clearpage
\subsection{S3 — tái dựng đoạn đường đã đi}\label{sub:4.3}
Đối thủ cần tái dựng các vị trí của đoạn đã đi. A là chuỗi đều; B ít nhánh; C quan sát thưa.

\noindent\begin{minipage}[t]{0.48\linewidth}\vspace{0pt}\includegraphics[width=\linewidth]{figures/case_S3_A.pdf}\end{minipage}\hfill\begin{minipage}[t]{0.49\linewidth}\vspace{0pt}\RaggedRight \textbf{S3.A / 12 bản ghi.} Đoạn di chuyển $\ge$60 giây qua $\ge$3 cạnh; gửi mỗi 20 giây.\par\smallskip \textbf{Bản ghi minh họa v2-r301-005.} u301\_\allowbreak{}00: 12 mẫu, chỉ số GPS FCD[361…581]\par\smallskip\textbf{Đọc sample.} 12 mẫu của u301\_\allowbreak{}00 cách nhau 20 giây. Chuỗi điểm cung cấp nhiều ràng buộc để nối tuyến; nét đứt là đường thật chỉ dùng chấm.\end{minipage}\par\vspace{3pt}
\noindent\begin{minipage}[t]{0.48\linewidth}\vspace{0pt}\includegraphics[width=\linewidth]{figures/case_S3_B.pdf}\end{minipage}\hfill\begin{minipage}[t]{0.49\linewidth}\vspace{0pt}\RaggedRight \textbf{S3.B / 11 bản ghi.} Ít nhất một nửa số cạnh được lấy mẫu chỉ có một hướng đi tiếp: hành lang ít lựa chọn.\par\smallskip \textbf{Bản ghi minh họa v2-r301-007.} u301\_\allowbreak{}00: 9 mẫu, chỉ số GPS FCD[91…251]\par\smallskip\textbf{Đọc sample.} 9 mẫu; 55,6\% cạnh được lấy mẫu chỉ có một hướng đi tiếp. Đường ít nhánh có thể khiến chuỗi dễ suy hơn dù từng vị trí được làm mờ.\end{minipage}\par\vspace{3pt}
\noindent\begin{minipage}[t]{0.48\linewidth}\vspace{0pt}\includegraphics[width=\linewidth]{figures/case_S3_C.pdf}\end{minipage}\hfill\begin{minipage}[t]{0.49\linewidth}\vspace{0pt}\RaggedRight \textbf{S3.C / 12 bản ghi.} Lấy mẫu đoạn di chuyển thưa hơn, mỗi 60 giây.\par\smallskip \textbf{Bản ghi minh họa v2-r301-006.} u301\_\allowbreak{}00: 4 mẫu, chỉ số GPS FCD[361…541]\par\smallskip\textbf{Đọc sample.} 4 mẫu của cùng đoạn trong A, cách nhau 60 giây. Khoảng trống lớn hơn nhưng đối thủ vẫn có thể dùng mạng đường để nối các khả năng.\end{minipage}\par\vspace{3pt}

\clearpage
\subsection{S9 — suy ra điểm xuất phát bị che}\label{sub:4.4}
Đối thủ cần suy điểm đầu bị giấu. A suy ngược một chuyến; B phân biệt hai nguồn nhập tuyến; C kết hợp chuyến lặp.

\noindent\begin{minipage}[t]{0.48\linewidth}\vspace{0pt}\includegraphics[width=\linewidth]{figures/case_S9_A.pdf}\end{minipage}\hfill\begin{minipage}[t]{0.49\linewidth}\vspace{0pt}\RaggedRight \textbf{S9.A / 12 bản ghi.} Cạnh xuất phát chỉ một lối ra. Che 60 giây đầu, lấy phần còn lại mỗi 20 giây.\par\smallskip \textbf{Bản ghi minh họa v2-r301-025.} u301\_\allowbreak{}00: 27 mẫu, chỉ số GPS FCD[60…580]\par\smallskip\textbf{Đọc sample.} Sao là FCD[0]; cửa sổ bắt đầu ở FCD[60] sau khi che 60 giây. Chỉ một lối ra có thể giúp suy ngược nguồn xuất phát từ phần còn công bố.\end{minipage}\par\vspace{3pt}
\noindent\begin{minipage}[t]{0.48\linewidth}\vspace{0pt}\includegraphics[width=\linewidth]{figures/case_S9_B.pdf}\end{minipage}\hfill\begin{minipage}[t]{0.49\linewidth}\vspace{0pt}\RaggedRight \textbf{S9.B / 11 bản ghi.} Hai chuyến khác điểm đầu nhưng nhập một đoạn chung; chỉ cấp từ chỗ nhập tuyến.\par\smallskip \textbf{Bản ghi minh họa v2-r301-026.} u301\_\allowbreak{}00: 28 mẫu, chỉ số GPS FCD[34…574]; u301\_\allowbreak{}04: 28 mẫu, chỉ số GPS FCD[19…559]\par\smallskip\textbf{Đọc sample.} Hai điểm đầu cách 281,76 m rồi nhập tuyến. Chỉ cấp dữ liệu sau nhập; hai sao là hai đáp án bị giấu mà đối thủ phải phân biệt.\end{minipage}\par\vspace{3pt}
\noindent\begin{minipage}[t]{0.48\linewidth}\vspace{0pt}\includegraphics[width=\linewidth]{figures/case_S9_C.pdf}\end{minipage}\hfill\begin{minipage}[t]{0.49\linewidth}\vspace{0pt}\RaggedRight \textbf{S9.C / 12 bản ghi.} Nhiều chuyến có cùng điểm xuất phát dự kiến; che 60 giây đầu từng chuyến.\par\smallskip \textbf{Bản ghi minh họa v2-r301-027.} u301\_\allowbreak{}00: 27 mẫu, chỉ số GPS FCD[60…580]; u301\_\allowbreak{}09: 27 mẫu, chỉ số GPS FCD[60…580]\par\smallskip\textbf{Đọc sample.} u301\_\allowbreak{}00/u301\_\allowbreak{}09 lặp cùng điểm đầu dự kiến; mỗi phiên che 60 giây đầu. Nhiều cửa sổ có thể bổ sung bằng chứng về cùng nguồn xuất phát.\end{minipage}\par\vspace{3pt}

\clearpage
\subsection{S10 — suy ra điểm kết thúc bị che}\label{sub:4.5}
Đối thủ cần suy điểm cuối của chuyến đã hoàn tất. A dùng một chuyến; B kết hợp nhiều chuyến cùng đích. Đây khác dự đoán đích tương lai ở S6.

\noindent\begin{minipage}[t]{0.58\linewidth}\vspace{0pt}\includegraphics[width=\linewidth]{figures/case_S10_A.pdf}\end{minipage}\hfill\begin{minipage}[t]{0.39\linewidth}\vspace{0pt}\RaggedRight \textbf{S10.A / 11 bản ghi.} Cạnh kết thúc chỉ một lối vào. Che 60 giây cuối của chuyến đã hoàn tất.\par\smallskip \textbf{Bản ghi minh họa v2-r302-025.} u302\_\allowbreak{}08: 38 mẫu, chỉ số GPS FCD[0…740]\par\smallskip\textbf{Đọc sample.} Sao là FCD cuối [817] của u302\_\allowbreak{}08; che 60 giây cuối. Đường chỉ một lối vào có thể giúp suy đoạn kết thúc từ phần trước còn thấy.\end{minipage}\par\vspace{3pt}
\noindent\begin{minipage}[t]{0.58\linewidth}\vspace{0pt}\includegraphics[width=\linewidth]{figures/case_S10_B.pdf}\end{minipage}\hfill\begin{minipage}[t]{0.39\linewidth}\vspace{0pt}\RaggedRight \textbf{S10.B / 12 bản ghi.} Nhiều chuyến cùng điểm kết thúc dự kiến; che 60 giây cuối từng chuyến.\par\smallskip \textbf{Bản ghi minh họa v2-r301-028.} u301\_\allowbreak{}00: 27 mẫu, chỉ số GPS FCD[0…520]; u301\_\allowbreak{}09: 27 mẫu, chỉ số GPS FCD[0…520]\par\smallskip\textbf{Đọc sample.} u301\_\allowbreak{}00/u301\_\allowbreak{}09 lặp cùng đích dự kiến và cùng che 60 giây cuối. Đối thủ có thể kết hợp nhiều phiên để khoanh đích, chưa đủ gọi đó là nơi làm việc.\end{minipage}\par\vspace{3pt}
S9/S10 chấm tọa độ đầu/cuối; chưa gán nhãn nhà/nơi làm việc. Xem bản đồ phóng to của 29 ca của khung đầy đủ ở sample\_\allowbreak{}maps.html; tọa độ và bản ghi ở data\_\allowbreak{}samples.json, số đếm ở scenario\_\allowbreak{}inventory.csv.


\clearpage
\section{Nguồn và khả năng truy nguyên}\label{sec:5}
Bảng rút gọn đọc artifacts/benchmarks/active\_\allowbreak{}scope\_\allowbreak{}ac\_\allowbreak{}v2/readout.json đã xác minh, không chạy thêm mô hình. method\_\allowbreak{}evidence.json giữ hash benchmark; preparation\_\allowbreak{}evidence.json giữ bảng guide. Hồ sơ metrics gốc nằm trong bản đầy đủ lưu ở archive/before\_\allowbreak{}concise\_\allowbreak{}2026-10-03/ và sources.json. Module concise\_\allowbreak{}presentation\_\allowbreak{}content.py tạo nội dung rút gọn; plot\_\allowbreak{}model\_\allowbreak{}architecture.py vẽ mô hình từ các component trong mã. scenario\_\allowbreak{}appendix.tex dùng chung cho hai tài liệu.

Ngày 03/10/2026 cập nhật cách trình bày, đối chiếu bao phủ scenario và làm rõ kiến trúc; không thay số liệu thực nghiệm. Tra data\_\allowbreak{}samples.json để phân biệt nhãn trình bày với mã thực nghiệm, đặc biệt S10.B (mã nguồn S10.C cũ).

\par\hypertarget{ref-R1}{}\textbf{[R1]} Yadav et al. (2024). Protecting Vehicle Location Privacy with Contextually-Driven Synthetic Location Generation (TransProtect). SIGSPATIAL. \href{https://arxiv.org/html/2409.09495v1}{Nguồn gốc}. Toàn văn tác giả; §5.1.4, Eq. 13.
\par\hypertarget{ref-R2}{}\textbf{[R2]} Liu, Peng, Zhou (2026). A Dummy-Based Location Privacy Protection Scheme with Semantic Correlation of Moving Paths. JKSUCIS 38:478. \href{https://link.springer.com/article/10.1007/s44443-026-00899-w}{Nguồn gốc}. Toàn văn NXB; §6.3–6.5.
\par\hypertarget{ref-R3}{}\textbf{[R3]} Liu, Hu, Zhou (2026). Location Privacy Protection in Continuous LBSs: Enhancing Anonymity via Fake Queries. JKSUCIS 38:53. \href{https://link.springer.com/article/10.1007/s44443-025-00438-z}{Nguồn gốc}. Toàn văn NXB; §6.3–6.6.
\par\hypertarget{ref-R4}{}\textbf{[R4]} Road Network-Aware Personalized Trajectory Protection with Differential Privacy under Spatiotemporal Correlations (2025), arXiv:2511.21020v1. \href{https://arxiv.org/html/2511.21020v1}{Nguồn gốc}. Bản tiền công bố; §VI, Eq. 26–27.
\par\hypertarget{ref-R5}{}\textbf{[R5]} Wang et al. (2025). CPCROK: A Communication-Efficient and Privacy-Preserving Scheme for Low-Density Vehicular Ad Hoc Networks. Future Internet 17(4):165. \href{https://uhra.herts.ac.uk/id/eprint/25705/1/futureinternet-17-00165.pdf}{Nguồn gốc}. Toàn văn lưu tại trường tác giả; §5.4.
\par\hypertarget{ref-R6}{}\textbf{[R6]} Yue, Li, Liu, Li (2025). DP-FETC: A differentially private trajectory publishing method based on feature extraction and trajectory correlation. ERA 33(11):6631–6651. \href{https://www.aimspress.com/aimspress-data/era/2025/11/PDF/era-33-11-293.pdf}{Nguồn gốc}. Toàn văn NXB; §4.4, §5.2, Eq. 19–21.
\par\hypertarget{ref-R7}{}\textbf{[R7]} Gu et al. (2025). Research on Joint Protection of LBS Location and Query Privacy in Internet of Vehicles Based on an Improved PIR Algorithm. Concurrency and Computation: Practice and Experience. \href{https://onlinelibrary.wiley.com/doi/abs/10.1002/cpe.70447}{Nguồn gốc}. Chỉ xác minh tóm tắt NXB; chưa đủ định nghĩa các metric.
\par\hypertarget{ref-R8}{}\textbf{[R8]} Buchholz et al. (2024). SoK: Can Trajectory Generation Combine Privacy and Utility? PoPETs 2024(3):75–93. \href{https://petsymposium.org/popets/2024/popets-2024-0068.pdf}{Nguồn gốc}. Toàn văn; khảo sát và hướng dẫn đánh giá, không phải một cơ chế đối chứng.
\par\hypertarget{ref-R9}{}\textbf{[R9]} Farahnakiyan, Esmaeilyfard, Javidan (2025). A proactive privacy-preserving framework for mobile trajectory sharing (PRISM). JNCA 242:104271. \href{https://www.sciencedirect.com/science/article/pii/S1084804525001687}{Nguồn gốc}. Tóm tắt/đoạn công khai của NXB; chưa xác minh công thức LPI/TDD/DC.
\par\hypertarget{ref-R10}{}\textbf{[R10]} Ioannou, Aktypi, Athanasopoulos (2026). From Options to Action: Evaluating Adoption of Privacy Features in Fitness-Tracking Platforms. CHI 2026. \href{https://elathan.github.io/papers/chi26.pdf}{Nguồn gốc}. Toàn văn tác giả; bảng 3, §6.4. Đo sử dụng tính năng, không đo sức chống tấn công.
\par\hypertarget{ref-B1}{}\textbf{[B1]} Niu et al. (2014). Achieving k-Anonymity in Privacy-Aware Location-Based Services (DLS). INFOCOM. \href{https://doi.org/10.1109/INFOCOM.2014.6848002}{Nguồn gốc}. Đối chứng nền; định nghĩa entropy theo hồ sơ khảo sát đã có trong repo.
\par\hypertarget{ref-B2}{}\textbf{[B2]} Shaham et al. (2021). Privacy Preservation in Location-Based Services: A Novel Metric and Attack Model (RDG). IEEE TMC 20(10). \href{https://doi.org/10.1109/TMC.2020.2993599}{Nguồn gốc}. Đối chứng nền; entropy/chuyển tiếp/Viterbi theo hồ sơ đã có trong repo.
\par\hypertarget{ref-B3}{}\textbf{[B3]} Li et al. (2024 issue; online 11/09/2023). AnotherMe: A Location Privacy Protection System Based on Online Virtual Trajectory Generation. IEEE TDSC 21(4):2552–2567. \href{https://ieeexplore.ieee.org/document/10246991/}{Nguồn gốc}. Metadata/tóm tắt và mã nguồn tác giả; chưa kiểm chứng đầy đủ mẫu số/đơn vị phép đo.
\par\hypertarget{ref-B4}{}\textbf{[B4]} Brauer et al. (2022). My home is my secret: concealing sensitive locations by context-aware trajectory truncation. IJGIS 36(12):2496–2524. \href{https://doi.org/10.1080/13658816.2022.2081694}{Nguồn gốc}. Toàn văn NXB và chương tác giả 2023; S-TT, k-site-unidentifiability. Ngoại lệ nền trực tiếp cho S9/S10.
\par\hypertarget{ref-R11}{}\textbf{[R11]} Forsch et al. (2023). Efficient Mining of Volunteered Trajectory Datasets, §3.2. \href{https://link.springer.com/chapter/10.1007/978-3-031-35374-1_3}{Nguồn gốc}. Toàn văn tác giả, online 09/12/2023. Trình bày lại S-TT 2022, không phải một phương pháp mới độc lập.
\par\hypertarget{ref-B5}{}\textbf{[B5]} Dhondt et al. (2022). A Run a Day Won’t Keep the Hacker Away: Inference Attacks on Endpoint Privacy Zones in Fitness Tracking Social Networks. CCS. \href{https://people.cs.kuleuven.be/~stijn.volckaert/papers/2022_CCS_Fitness_Tracking.pdf}{Nguồn gốc}. Bản tác giả: tấn công EPZ và sáu countermeasures; không gắn nhãn một model bảo vệ đã thắng benchmark.
\par\hypertarget{ref-B6}{}\textbf{[B6]} Dehaene et al. (2023). Masking Location Streams in the Presence of Colluding Service Providers. EuroS\&P Workshops. \href{https://lirias.kuleuven.be/bitstream/20.500.12942/720391/2/Data_Stream_Obfuscator_camera_ready.pdf}{Nguồn gốc}. Bản tác giả: Data Stream Obfuscator, privacy zones/timeslots, mapping cache. Tháng 7/2023, ngoài cửa sổ ba năm.
\par\hypertarget{ref-R12}{}\textbf{[R12]} Atmaca et al. (2024). A Privacy-Preserving Querying Mechanism with High Utility for Electric Vehicles. IEEE OJVT. \href{https://wrap.warwick.ac.uk/id/eprint/183198/1/WRAP-privacy-preserving-querying-mechanism-high-utility-electric-vehicles-2024.pdf}{Nguồn gốc}. Bản tác giả: AGeoI + dummy cho truy vấn trạm sạc; không xác minh benchmark điểm đầu/cuối.
\end{document}
